\section{Conclusion}
\label{sec:conclusion}

Retirement policies change behavior over the entire working life, their effects depend on health, and their fiscal consequences spread across the tax and transfer system. This paper measures the marginal value of public funds of retirement reforms in an estimated life-cycle model with these features. For each reform, it computes the ex ante willingness to pay of workers who do not yet know their health and wage histories and the effect of the reform on Social Security outlays and on the net cost of income and payroll taxes, disability insurance, and means-tested transfers.

For the cohorts we study, one distinction organizes the ranking of reforms: whether a reform rewards later claiming or lowers benefits at every claiming age. Raising the delayed retirement credit, or steepening the adjustment of benefits around the normal retirement age, pays for itself. These reforms raise welfare and lower Social Security outlays; for the increases in the credit, outlays fall even though they would rise at unchanged behavior, so the claiming responses more than finance the higher credit. Raising the normal retirement age saves money by lowering benefits for every claimant and reduces welfare. Among expansions of benefits, an increase in the bottom replacement rate of the benefit formula delivers the most welfare per dollar; the smallest such increase is approximately budget neutral. Expansions of Medicare coverage after age 65 deliver the least, and cuts in old-age income taxes little more.

These results suggest an order of operations for closing a financing gap under the rules we study. Reforms that reward delay come first, because they raise revenue and welfare together. To a first-order approximation, an increase in the normal retirement age costs less welfare per dollar saved than higher income taxes on older workers, but in our grid it is dominated by rewards for delay whenever the required saving is moderate. The two changes to claiming incentives in the 1983 Amendments illustrate the distinction: evaluated for the 1920--1940 cohorts, the end point of the increase in the delayed retirement credit pays for itself, while the end point of the increase in the normal retirement age is a benefit cut with a low MVPF.

Four extensions would show how far these conclusions extend beyond the cohorts we study. The first is to evaluate rewards for delay at current-law parameters, with a delayed retirement credit of 8 percent and a normal retirement age of 67, and with survival beyond age 85. The second is to weight welfare gains by health and lifetime earnings rather than averaging them ex ante. The third is to allow wages and the government budget to respond in general equilibrium. The fourth is to let statutory retirement ages act as reference points. The logic behind our results is more general than the model: when the cost of continued work differs across workers, a reward for later claiming pays only the workers who choose to delay, so it can lower outlays when delays are cheap for the program, whereas an increase in the normal retirement age lowers the benefits of all claimants, including those for whom working longer is most costly.
